%------------------------------------------------------------------------------%
\section{Teste da Comparação}

Nessa seção apresentamos um resultado que nos permite comparar séries conhecidas
com séries que queremos analisar, o Teste da Comparação. Existem duas versões
desse teste o \emph{Teste da Comparação} apresentado no teorema a seguir
e o \emph{Teste da Comparaçao no Limite} apresentado no Teorema~\ref{teo:teste-comparacao-limite}.

%------------------------------------------------------------------------------%
\begin{teorema}{Teste da Comparação}{teo:teste-comparacao}
  Sejam%
  \index{Teste!comparação}
  \[
    \sum_{n=N}^\infty a_n
    \qquad\text{ e }\qquad
    \sum_{n=N}^\infty b_n
  \]
  séries tais que $a_n\geq 0$ e $b_n\geq 0$. Suponha que para algum~$N$,
  \[
    a_n\leq b_n, \qquad \forall  n\geq N
  \]
  então,
  \begin{enumerate}
    \item se \(\sum_{n=N}^\infty b_n\) converge, \tabto{45mm} então \(\sum_{n=N}^\infty a_n\) converge;
    \item se \(\sum_{n=N}^\infty a_n\) diverge,  \tabto{45mm} então \(\sum_{n=N}^\infty b_n\) diverge.
  \end{enumerate}
\end{teorema}
%------------------------------------------------------------------------------%
\begin{proof}
  \emph{Parte 1}

  Como
  \[
    a_n\leq b_n, \quad \forall  n\geq N
  \]
  temos que
  \[
    \sum_{k=N}^n a_k    \;\leq\;
    \sum_{k=N}^n b_k    \qquad \forall n\geq N      \,.
  \]
  Somando $b_1+b_2+\cdots+ b_{N-1}$ de ambos os lados, obtemos
  \[
    b_1 + b_2 + \cdots + b_{N-1} + \sum_{k=N}^na_k
    \;\leq\;
    \sum_{k=1}^nb_k=B_n \qquad \forall n\geq N         \,.
  \]
  O termo $B_n$, do lado direito da desigualdade acima, é a sequência de
  somas parciais da série $\sum b_n$, como essa série converge e
  $b_n\geq 0$ (ou seja a $B_n$ é crescente), segue que
  \[
    0 \;\leq\; b_1 + b_2 + \cdots + b_{N-1} + \sum_{k=N}^na_k
      \;\leq\; \sum_{k=1}^\infty b_k
      = B
      \qquad \forall n\geq N          \,,
  \]
  ou seja, a sequência
  \[
    b_1 + b_2 + \cdots + b_{N-1} + \sum_{k=N}^na_k
  \]
  é limitada e crescente,
  pois $a_n\geq 0$, portanto esta é uma sequência convergente, ou seja, a série
  \[
    b_1+b_2+\cdots+ b_{N-1}+\sum_{n=N}^\infty a_n
  \]
  converge. Mas como já discutimos
  anteriormente, adicionar e retirar um número finito determos em uma série não muda sua
  natureza, de modo que podemos retirar os termos $b_1+b_2+\cdots+ b_{N-1}$ e adicionar os
  termo $a_1+a_2+\cdots+ a_{N-1}$ e concluir que a série $\sum a_n$ é convergente.

\emph{Parte 2}

  Suponha por absurdo que $\sum b_n$ converge, então, pelo item anterior,
  segue que $\sum a_n$ converge. O que é um absurdo, pois por hipótese $\sum a_n$
  diverge. Segue que $\sum b_n$ diverge.
\end{proof}
%------------------------------------------------------------------------------%

%------------------------------------------------------------------------------%
\begin{example}
  Mostre que a série \(\sum_{n=1}^\infty \frac{n-1}{n^4+2}\) converge.
  \solution
  De fato, veja que
  \begin{align*}
    \frac{n-1}{n^4+2} & =    \frac{n}{n^4+2} - \frac{1}{n^4+2} \\[3mm]
                      & \leq \frac{n}{n^4+2}                   \\[3mm]
                      & =    \frac{1}{n^3 + \sfrac{2}{n}}      \\[3mm]
                      & \leq \frac{1}{n^3}                     \,.
  \end{align*}
  Como a série
  \[
    \sum_{n=1}^\infty \frac{1}{n^3}
  \]
  é uma série $p$ com $p=3>1$, ela é convergente.
  Assim, pelo teste de comparação a série
  \[
    \sum_{n=1}^\infty \frac{n-1}{n^4+2}
  \]
  também converge.
\end{example}
%------------------------------------------------------------------------------%

%------------------------------------------------------------------------------%
\begin{example}
  Verifique a convergência da série
  \(
    \sum_{n=1}^\infty\frac{3^n}{2^n-1}
  \).
  \solution
  A série
  \[
    \sum_{n=1}^\infty\frac{3^n}{2^n-1}
  \]
  diverge, pois
  \[
    \frac{3^n}{2^n-1}\geq \left(\frac{3}{2}\right)^n
  \]
  e a série geométrica
  \[
    \sum_{n=1}^\infty\left(\frac{3}{2}\right)^n
  \]
  diverge, pois
  \(
    \abs{r}=\sfrac{3}{2}>1
  \).
\end{example}
%------------------------------------------------------------------------------%

Apresentamos agora a segunda versão do teste da comparação. nessa versão substituímos
a necessidade de verificar se $a_n\leq b_n$ ou $a_n\geq b_n$ pelo cálculo de um
limite. Ao aplicarmos o teste em uma série específica podemos escolher a versão
que necessitar de cálculos mais simples.

%------------------------------------------------------------------------------%
\begin{teorema}{Teste da Comparação no Limite}{teo:teste-comparacao-limite}
  Sejam\index{Teste!comparação no limite}
  $a_n>0$ e $b_n>0$, para todo $n\geq N$, $N\in \N$.
  \begin{enumerate}
    \item
      Se
      \(
        \lim\frac{a_n}{b_n}=c>0
      \;\)
      então
      \(\;
        \sum_{n=1}^\infty a_n
      \)
      converge se, e somente se,
      \(
        \sum_{n=1}^\infty b_n
      \)
      converge.
    \item
      Se
      \(
        \lim\frac{a_n}{b_n}=0
      \;\) e
      \(\;
        \sum_{n=1}^\infty b_n
      \)
      converge, então
      \(
        \sum_{n=1}^\infty a_n
      \)
      converge.
    \item
      Se
      \(
        \lim\frac{a_n}{b_n}=\infty
      \;\) e
      \(\;
        \sum_{n=1}^\infty b_n
      \)
      diverge, então
      \(
        \sum_{n=1}^\infty a_n
      \)
      diverge.
  \end{enumerate}
\end{teorema}
%------------------------------------------------------------------------------%
\begin{proof}
  %----------------------------------------------------------------------------%
  \emph{Parte 1}

  A demonstração da Parte 1 está no Teorema 11 na Seção 10.4 no livro do Thomas~\cite{thomas_2}.

  %----------------------------------------------------------------------------%
  \emph{Parte 2}

  Para todo $\varepsilon>0$ existe $N_1\in\N$,  $N_1\geq N$, tal que
  \[
    \frac{a_n}{b_n} =    \abs{\frac{a_n}{b_n}-0}
                    \leq \varepsilon
    \qquad \forall n\geq N_1
  \]
  ou seja,
  \[
    a_n \leq \varepsilon b_n \qquad \forall  n\geq N_1          \,.
  \]
  Assim, pelo teorema anterior, se $\sum \varepsilon b_n$
  converge, então $\sum a_n$
  também converge, mas é claro que $\sum \varepsilon b_n$
  converge, pois $\sum b_n$
  converge e $\varepsilon$ é uma constante.

  %----------------------------------------------------------------------------%
  \emph{Parte 3}

  Para todo $M>0$ dado, existe $N_1\in\N$, $N_1\geq N$, tal que
  \[
    \frac{a_n}{b_n}>M \qquad \forall n> N_1   \,,
  \]
  ou seja,
  \[
    a_n > M b_n \qquad \forall n > N_1     \,.
  \]
  Como $\sum b_n $ diverge e $M$ é constante, segue que $\sum M b_n$ também diverge.
  Pelo teorema anterior, como $\sum M b_n$ diverge, segue que $\sum a_n$ também diverge.
\end{proof}
%------------------------------------------------------------------------------%

%------------------------------------------------------------------------------%
\begin{example}
  Verificar a convergência da série
  \(
    \sum_{n=1}^\infty \frac{2^n}{3+4^n}
  \).
  \solution
  A série
  \[
    \sum_{n=1}^\infty \frac{2^n}{3+4^n}
  \]
  converge, pois tomando
  \[
    b_n=\frac{2^n}{4^n}=\frac{1}{2^n}
  \]
  temos
  \[
    \lim \frac{a_n}{b_n} = \lim \frac{4^n}{3+4^n}
                         = \lim \frac{1}{\sfrac{3}{4^n}+1}
                         = 1                  \,.
  \]
  Como a série
  \[
    \sum \frac{1}{2^n}
  \]
  converge, pois é uma série geométrica com $\abs{r}=\sfrac{1}{2}<1$,
  segue, usando o item 1 do teorema, que a série
  \[
    \sum_{n=1}^\infty \frac{2^n}{3+4^n}
  \]
  também converge.
\end{example}
%------------------------------------------------------------------------------%

%------------------------------------------------------------------------------%
\begin{example}
  Use o fato que $\ln n$ cresce mais lentamente que $n^c$ para qualquer
  constante positiva $c$ para provar a convergência da série
  \[
    \sum_{n=1}^\infty \frac{(\ln n)^2}{n^3}      \,.
  \]
  \solution
  Para verificarmos que $\ln n$ cresce mais lentamente que $n^c$ para qualquer
  constante positiva $c$ usamos a regra de L'Hôpital para mostrar que
  \[
    \lim_{x\to \infty}\frac{\ln x}{x^c} = 0   \,.
  \]
  Usando esse fato podemos escrever que para qualquer $c>0$ vale a desigualdade
  \[
    \frac{\left(\ln n\right)^2}{n^3} < \frac{n^{2c}}{n^3} = \frac{1}{n^{3-2c}}     \,.
  \]
  Queremos agora escolher a constante $c$ tal que a expressão do lado direito
  seja o termo geral de uma série convergente.
  Observando que esses são os termos de uma série-$p$ com $p=3-2c$, precisamos
  impor que $3-2c>1$ ou seja $c<1$, escolhemos então \(c=\sfrac{1}{2}\).
  Dessa forma construímos uma série convergente com termos
  \[
     b_n = \frac{1}{n^{3-2c}}
         = \frac{1}{n^2}     \,.
  \]
  Podemos aplicar o Teste da Comparação no Limite
  \begin{align*}
    \lim \frac{a_n}{b_n} & = \lim \frac{(\ln n)^2} {n} \\[3mm]
                         & =2\lim \frac{\ln n}{n}      \\[3mm]
                         & =2\lim \frac{ 1}{n}=0       \,.
  \end{align*}
  Pelo item $2$ do Teorema~\ref{teo:teste-comparacao-limite},
  como
  \[
    \sum_{n=1}^\infty \frac{1}{n^2}
  \]
  converge, segue que a série em questão também converge.
\end{example}
%------------------------------------------------------------------------------%

%------------------------------------------------------------------------------%
\begin{example}
    Verifique que série
    \(
      \sum_{n=1}^\infty \left(\frac{2n+3}{5n+4}\right)^n
    \)
    diverge.

    \solution
    Com efeito, comparando o termo geral
    \(
      a_n=\left(\frac{2n+3}{5n+4}\right)^n
    \)
    com
    \(
      b_n=\left(\frac{2}{5} +\frac{3}{5n}\right)^n
    \),
    temos
    \begin{align*}
      \lim \frac{a_n}{b_n} & = \lim \left( \frac{2n+3}{5n+4} \; \frac{5n}{2n+3}\right)^n \\
                           & = \lim \left( \frac{5n}{5n+4}\right)^n                      \\
                           & = \lim \left( \frac{1}{1+\sfrac{4}{5n}}\right)^n            \\
                           & = \frac{1}{e^{\sfrac{4}{5}}}                                \,.
    \end{align*}
    Como
    \(
      \frac{1}{e^{\sfrac{4}{5}}} > 0
    \),
    pelo primeiro item do teorema, segue que o comportamento
    de $\sum a_n$ é o mesmo de $\sum b_n$, mas $\sum b_n$ diverge, pois
    \[
      b_n=\left(\frac{2}{5} +\frac{3}{5n}\right)^n\to e^{\sfrac{6}{25}}\neq 0  \,.
    \]
\end{example}
%------------------------------------------------------------------------------%

Para verificar as passagens desse exemplo, é possível utilizar a expressão
\[
  \lim \left( \frac{1}{1+\sfrac{4}{5n}}\right)^n =
  \exp\left[ \,\lim\ln \left( \dfrac{1}{1+\sfrac{4}{5n}}\right)^n \,\right]    \,.
\]

%------------------------------------------------------------------------------%
